\mmedAbsLabel{目的：}
本公开示例以“基于公开重症监护数据的脓毒症早期风险分层流程”为载体，演示南方医科大学博士后研究报告在封面、题名页、摘要、目录、正文、参考文献、附录与结尾部分上的完整组织方式。

\mmedAbsLabel{方法：}
示例采用“问题定义、变量整理、流程分层、结果展示、局限分析”的五段式写法，并用公开场景描述、占位图表和标准化前后置页结构模拟正式研究报告的基本排版需求。模板层面继续复用 \texttt{bensz-thesis} 的成熟页眉、目录、图表、参考文献与构建脚本，并在包级注册独立的 \texttt{thesis-smu-postdoc} template/profile/style，确保单项目分发时仍保持完整模板身份。

\mmedAbsLabel{结果：}
最终得到一套可直接编译的南方医科大学博士后研究报告项目脚手架。其封面吸收了《博士后研究报告编写规则》中“分类号 / UDC / 密级 / 编号 / 工作完成日期 / 报告提交日期”的核心字段，同时在题名页中补齐博士后姓名、流动站、专业、起止时间和单位信息。正文示例覆盖章节、表格、示意图、附录、个人简历、成果清单与永久通信地址等关键结构。

\mmedAbsLabel{结论：}
该模板证明，在继续复用现有 thesis 公共基础设施的前提下，博士后研究报告也可以以独立 template ID 的方式稳定交付，同时保持单项目与完整仓库两种使用方式的一致性。

\mmedAbsLabel{关键词：}
博士后研究报告；南方医科大学；LaTeX 模板；公开示例；结构化排版
